\chapter{Methodology}
\label{chap:methodology}

Questo capitolo descrive la metodologia adottata per costruire la base
sperimentale della tesi e per predisporre la valutazione della robustezza
operativa di sistemi \gls{ai}-based per la classificazione automatizzata
di immagini in ambito forense. In coerenza con il gap individuato nel
Capitolo~\ref{chap:state-of-the-art}, relativo alla limitata disponibilità
di protocolli riproducibili per validare modelli e strumenti \gls{ai}-based
in workflow forensi realistici, il capitolo non si limita a descrivere
singoli componenti tecnici, ma formalizza un disegno sperimentale integrato.
Tale disegno collega dataset, modelli, perturbazioni adversariali e
anti-forensi, strumenti forensi, metriche, approccio di \gls{xai} e controlli
di tracciabilità all'interno del framework \gls{fairlab}, mantenendo una
distinzione netta tra metodologia sperimentale e analisi dei risultati.
Il capitolo è organizzato in tredici sezioni: le prime sei
(\S\S\,\ref{sec:methodology-fairlab-overview}--\ref{sec:methodology-split-generation})
descrivono le componenti della pipeline dataset; le successive sette
(\S\S\,\ref{sec:methodology-models-classification}--\ref{sec:methodology-traceability-controls})
introducono il paradigma di classificazione, la selezione e la generazione
delle perturbazioni, il \emph{forensic evaluation bundle}, le metriche
di valutazione, l'approccio di \gls{xai} e i controlli di tracciabilità
e riproducibilità.

La trattazione del presente capitolo si concentra sulle componenti
metodologiche già consolidate nella pipeline dataset: acquisizione delle
sorgenti, preparazione tecnica del pool di immagini, generazione del manifest
di revisione, selezione manuale \textit{human-in-the-loop}, congelamento del
dataset finale, costruzione del subset binario per gli esperimenti adversariali
e generazione degli split puliti con set \gls{ood} separato. Le fasi
successive, relative alla generazione delle perturbazioni, alla costruzione
del \textit{forensic evaluation bundle} e alla valutazione mediante strumenti
forensi commerciali, sono richiamate soltanto come raccordo metodologico e
saranno sviluppate nei passaggi sperimentali successivi.


% ─────────────────────────────────────────────────────────────────────────────
\section{Overview of the FAIR-Lab methodology}
\label{sec:methodology-fairlab-overview}
% ─────────────────────────────────────────────────────────────────────────────

La metodologia della tesi è organizzata attorno al framework \gls{fairlab},
inteso come pipeline sperimentale per la valutazione della robustezza operativa
di sistemi \gls{ai}-based in scenari di \gls{df}. Il framework nasce
dall'esigenza di superare una valutazione puramente algoritmica della
classificazione di immagini, collegando le prestazioni dei modelli alla loro
possibile funzione operativa nei workflow di triage e prioritizzazione delle
evidenze digitali. In tale prospettiva, la robustezza non viene considerata
soltanto come proprietà tecnica del classificatore, ma come requisito
metodologico per l'impiego prudente dell'automazione in un contesto
investigativo.

Il framework integra due livelli complementari. Il primo è una pipeline
accademica, nella quale dataset, modelli, perturbazioni, metriche e approcci
di \gls{xai} vengono definiti in modo controllato e comparabile. Il secondo è
una pipeline operativo-forense, nella quale gli artefatti prodotti dalla
pipeline accademica vengono predisposti per essere analizzati da strumenti
professionali di \gls{df}, secondo una logica più vicina alle condizioni
applicative reali. Questa distinzione consente di mantenere separati il
controllo sperimentale proprio della ricerca accademica e la valutazione
black-box o quasi black-box richiesta dall'analisi di strumenti commerciali,
spesso caratterizzati da moduli proprietari e non pienamente ispezionabili
\cite{sanna2024preliminary,sanna2025pixelperfect}.

Nel suo stato attuale, la pipeline è stata consolidata fino alla costruzione
degli artefatti dataset, alla generazione degli split puliti e alla produzione
delle principali famiglie di perturbazioni adversariali e anti-forensi sul
subset binario. In particolare, sono stati definiti e prodotti il dataset
finale congelato, il subset binario destinato agli esperimenti adversariali
e anti-forensi, gli split puliti per la valutazione baseline e il set
\gls{ood} separato. Tali artefatti costituiscono la base stabile sulla quale
verranno eseguite le successive fasi di perturbazione, valutazione dei modelli,
analisi con strumenti forensi e \gls{xai}. Questa separazione tra artefatti
già congelati e fasi sperimentali successive è essenziale per evitare che
modifiche future ai dati alterino retroattivamente il protocollo di
valutazione.


% ─────────────────────────────────────────────────────────────────────────────
\section{Research design and experimental assumptions}
\label{sec:methodology-research-design}
% ─────────────────────────────────────────────────────────────────────────────

Il disegno sperimentale assume come scenario di riferimento un'attività di
triage visuale in ambito forense, nella quale un sistema automatico supporta
l'analista nell'individuazione preliminare di immagini potenzialmente
rilevanti. Il caso di studio riguarda la classificazione di immagini in tre
categorie metodologiche: \texttt{weapon}, \texttt{non\_weapon} e \texttt{ood}.
La prima categoria rappresenta la classe investigativamente rilevante nel task
considerato; la seconda raccoglie esempi negativi realistici e coerenti con il
dominio; la terza comprende casi borderline, anomali o fuori distribuzione che
non dovrebbero essere forzati in modo acritico nelle due classi binarie
principali.

La scelta di includere una classe \gls{ood} esplicita risponde a un'esigenza
metodologica specifica. In un contesto operativo reale, i sistemi di
classificazione non ricevono soltanto immagini pulite e semanticamente allineate
alle classi previste, ma anche contenuti ambigui, degradati, sintetici, non
pertinenti o appartenenti a distribuzioni differenti. La letteratura sulla
\gls{ood} detection ha mostrato che i classificatori possono produrre predizioni
ad alta confidenza anche su input estranei al dominio di addestramento,
rendendo necessario distinguere gli errori su classi note dalla gestione di
campioni che non dovrebbero essere classificati come positivi o negativi ordinari
\cite{hendrycks2017baseline,fort2021exploring}.

La metodologia distingue inoltre tra classificazione automatica, robustezza e
validazione operativa. La classificazione automatica riguarda la capacità del
sistema di assegnare una classe a un'immagine. La robustezza riguarda la
stabilità di tale comportamento quando l'input viene degradato, alterato o
manipolato. La validazione operativa riguarda invece l'impatto dell'output
automatico sul workflow forense, in particolare rispetto al rischio di falsi
negativi su contenuti rilevanti, falsi positivi che aumentano il carico
analitico e classificazioni improprie di campioni \gls{ood}. Questa distinzione
permette di evitare che l'accuratezza nominale su dati puliti venga interpretata
come garanzia sufficiente di affidabilità in scenari investigativi.

Il protocollo adotta infine un'impostazione \textit{human-in-the-loop}.
La componente manuale non viene trattata come elemento esterno o residuale,
ma come parte documentata della metodologia. In ambito forense, la costruzione
del dataset non può dipendere esclusivamente da etichette automatiche o da
directory originarie delle sorgenti: è necessario che i criteri semantici di
inclusione, esclusione e assegnazione alle classi siano espliciti, tracciabili
e coerenti con il task investigativo.


% ─────────────────────────────────────────────────────────────────────────────
\section{Dataset construction and source consolidation}
\label{sec:methodology-dataset-construction}
% ─────────────────────────────────────────────────────────────────────────────

La costruzione del dataset è stata articolata in una sequenza di fasi distinte,
progettate per separare la raccolta delle sorgenti, la preparazione tecnica dei
file, la revisione semantica e il congelamento degli artefatti finali. Questa
separazione è metodologicamente rilevante perché consente di distinguere ciò
che appartiene alla qualità tecnica dell'immagine da ciò che riguarda la sua
interpretazione semantica rispetto al task forense.

La fase iniziale ha previsto l'organizzazione di cinque sorgenti eterogenee
all'interno della directory \texttt{datasets/raw/}. Le sorgenti comprendono
dataset pubblici, materiale derivato da repository specializzate, raccolte
ottenute mediante scraping controllato, contenuti acquisiti secondo una logica
\gls{osint} e dati provenienti da fonti deep web-oriented. L'obiettivo non era
costruire un benchmark visivamente omogeneo, ma un pool iniziale sufficientemente
vario da rappresentare alcune delle condizioni di eterogeneità che caratterizzano
le indagini digitali contemporanee.

\begin{table}[H]
\centering
\caption{Sorgenti grezze utilizzate nella pipeline dataset.}
\label{tab:methodology-raw-sources}
\begin{tabularx}{\textwidth}{@{}lX@{}}
\toprule
\textbf{Directory sorgente} & \textbf{Ruolo metodologico} \\
\midrule
\texttt{01\_kaggle\_weapon}    & Dataset pubblico orientato al riconoscimento
                                 di armi, utile come sorgente strutturata
                                 iniziale. \\
\texttt{02\_deepfirearm}       & Sorgente specializzata per immagini di armi
                                 da fuoco, utile a includere esempi positivi
                                 più controllati. \\
\texttt{03\_google\_scraped}   & Raccolta web ottenuta tramite scraping,
                                 finalizzata a introdurre variabilità semantica
                                 e visiva. \\
\texttt{04\_telegram\_youtube} & Raccolta \gls{osint}-oriented da fonti social
                                 e video, utile a simulare contenuti meno
                                 strutturati. \\
\texttt{05\_deepweb}           & Sorgente deep web-oriented, impiegata per
                                 rappresentare materiale eterogeneo e non
                                 convenzionale. \\
\bottomrule
\end{tabularx}
\end{table}

A partire dalle sorgenti grezze, lo script
\texttt{datasets/scripts/prepared/08\_build\_prepared\_dataset.py} costruisce
il pool preparato in \texttt{datasets/prepared/final\_pool/}. Questa fase non
assegna etichette semantiche definitive, ma produce una base tecnica validata
attraverso controlli di leggibilità, filtraggio dimensionale, calcolo degli hash
e rimozione dei duplicati esatti. In particolare, ogni immagine candidata viene
verificata come file immagine valido, sottoposta a soglia minima di dimensione,
associata a un valore SHA256 e copiata nel pool preparato solo se supera i
controlli tecnici.

Il risultato principale della fase di preparazione è il file
\texttt{metadata.csv}, che registra per ciascun campione informazioni quali
\texttt{image\_id}, nome del file preparato, sorgente di provenienza, percorso
relativo originario, hash SHA256, dimensioni, estensione e dimensione del file.
La disponibilità di tali metadati consente di mantenere un collegamento stabile
tra la sorgente originaria e la versione normalizzata inclusa nel pool
sperimentale. Gli output di reportistica, tra cui \texttt{invalid\_images.csv},
\texttt{duplicates\_discarded.csv} e \texttt{prepared\_build\_summary.json},
documentano invece le esclusioni tecniche e i duplicati rimossi.

Questa fase ha una funzione esclusivamente tecnica. La pipeline evita
deliberatamente di inferire la classe finale di un'immagine dalla cartella
sorgente, dal nome del file o da eventuali pre-label automatici. Tale scelta
riduce il rischio di propagare bias o errori delle sorgenti originarie nella
selezione finale, mantenendo la distinzione tra provenienza del campione e
assegnazione semantica verificata.

Il manifest operativo di revisione comprendeva complessivamente 15\,321 immagini
candidate. La distribuzione tra le sorgenti era fortemente asimmetrica: la
sorgente \texttt{02\_deepfirearm} rappresentava da sola la componente
numericamente prevalente del pool, mentre le restanti sorgenti contribuivano
volumi significativamente inferiori. Tale asimmetria è stata trattata come una
caratteristica del pool candidato e non come un criterio di campionamento
finale: la fase \textit{human-in-the-loop} ha applicato criteri semantici e di
bilanciamento indipendenti dal volume originario delle singole sorgenti.
Di conseguenza, la sorgente \texttt{02\_deepfirearm}, pur essendo dominante nel
pool iniziale, contribuisce al dataset congelato con 100 immagini
\texttt{weapon}, pari alla stessa quota positiva assegnata alle altre sorgenti.
La distribuzione finale per sorgente e classe è riportata nella
Tabella~\ref{tab:final-dataset-source-distribution} della
Sezione~\ref{sec:methodology-human-review}.


% ─────────────────────────────────────────────────────────────────────────────
\section{Review manifest generation}
\label{sec:methodology-review-manifest}
% ─────────────────────────────────────────────────────────────────────────────

Dopo la costruzione del pool preparato, la pipeline genera un manifest di
revisione completo mediante lo script
\texttt{datasets/scripts/prepared/09\_generate\_review\_manifest\_full.py}.
Questo passaggio costituisce il raccordo tra la preparazione tecnica del dataset
e la successiva selezione manuale. Il file prodotto,
\texttt{datasets/prepared/manifests/review\_manifest\_full.csv}, conserva le
informazioni tecniche provenienti da \texttt{metadata.csv} e introduce colonne
dedicate alla revisione semantica.

Il manifest di revisione è progettato per supportare una pipeline controllata e
ispezionabile. Da un lato, mantiene campi di identità e provenienza, quali
\texttt{image\_id}, \texttt{sha256}, \texttt{relative\_path},
\texttt{source\_dataset} e \texttt{source\_relative\_path}. Dall'altro,
introduce campi destinati allo stato della revisione, all'etichetta finale, alle
note del revisore, all'identificativo del revisore e alla gestione dei casi
\gls{ood} o esclusi. Il manifest predispone inoltre colonne per eventuali
pre-label ausiliari, che tuttavia non sostituiscono la decisione manuale del
revisore: l'unico criterio di assegnazione della classe finale è il giudizio
semantico esplicito formulato durante la revisione \textit{human-in-the-loop}.
In questo modo, il manifest diventa l'artefatto centrale per documentare il
passaggio da pool tecnico a dataset semanticamente validato.

La scelta di introdurre i campi di revisione in un file separato rispetto al
metadata tecnico consente di preservare la distinzione tra due livelli della
pipeline. Il primo livello riguarda l'esistenza, la leggibilità e l'unicità
tecnica dei file. Il secondo riguarda invece la loro pertinenza rispetto al task
forense. Questa separazione è particolarmente importante perché un'immagine
tecnicamente valida può essere semanticamente inadatta alla classe \texttt{weapon}
o \texttt{non\_weapon}, mentre un'immagine borderline può essere utile come
campione \texttt{ood} per valutare l'affidabilità del sistema in condizioni
aperte.


% ─────────────────────────────────────────────────────────────────────────────
\section{Human-in-the-loop review and final dataset freezing}
\label{sec:methodology-human-review}
% ─────────────────────────────────────────────────────────────────────────────

La selezione finale del dataset è stata eseguita mediante un protocollo
\textit{human-in-the-loop}, implementato nello script
\texttt{datasets/scripts/final/10\_manual\_selection\_protocol\_reviewer.py}.
Tale protocollo consente al revisore di ispezionare manualmente le immagini del
pool preparato, assegnando una delle etichette finali previste: \texttt{weapon},
\texttt{non\_weapon}, \texttt{ood} o \texttt{exclude}. Le decisioni del
revisore vengono salvate in un database operativo e registrate in log dedicati,
in modo da rendere documentabile il processo di selezione.

I criteri semantici adottati sono stati definiti in modo restrittivo.
La classe \texttt{weapon} include soltanto armi individuali realistiche,
chiaramente visibili e coerenti con il task di classificazione; la valutazione
si basa sull'apparenza visiva dell'oggetto nell'immagine e non sulla possibilità
di accertarne la natura materiale, salvo i casi in cui elementi visivi evidenti
ne suggeriscano la natura di replica, giocattolo, modello airsoft o contenuto
sintetico, che vengono ricondotti alla classe \texttt{ood}. La classe
\texttt{non\_weapon} comprende immagini realistiche prive di armi, ma plausibili
come esempi negativi nel dominio investigativo considerato. La classe \texttt{ood}
raccoglie invece contenuti borderline o fuori distribuzione, come coltelli,
spade, repliche, airsoft, immagini sintetiche, videogiochi, mezzi militari
pesanti, esplosioni, scene di guerra, immagini fortemente degradate o contenuti
anomali rispetto al dominio di interesse. La classe \texttt{exclude} è riservata
a file non utilizzabili o non coerenti con il dataset finale.

Il risultato della revisione è il dataset frozen ufficiale, registrato nel
manifest \texttt{datasets/final/manifests/manual\_selection\_final\_1500.csv}.
Il dataset finale contiene 1500 immagini bilanciate nelle tre classi principali,
come indicato nella Tabella~\ref{tab:methodology-final-dataset}. La
distribuzione per sorgente e classe, riportata nella
Tabella~\ref{tab:final-dataset-source-distribution}, mostra come la selezione
manuale abbia compensato la sproporzione del pool candidato: ogni sorgente
contribuisce alla classe \texttt{weapon} con 100 immagini, mentre la
distribuzione delle classi \texttt{non\_weapon} e \texttt{ood} riflette le
caratteristiche semantiche delle singole sorgenti.

\begin{table}[H]
\centering
\caption{Distribuzione del dataset frozen ufficiale.}
\label{tab:methodology-final-dataset}
\begin{tabularx}{0.75\textwidth}{@{}l>{\raggedleft\arraybackslash}X@{}}
\toprule
\textbf{Classe} & \textbf{Numero di immagini} \\
\midrule
\texttt{weapon}      & 500 \\
\texttt{non\_weapon} & 500 \\
\texttt{ood}         & 500 \\
\midrule
\textbf{Totale}      & \textbf{1500} \\
\bottomrule
\end{tabularx}
\end{table}

\begin{table}[H]
\centering
\caption{Distribuzione del dataset finale congelato per sorgente e classe.}
\label{tab:final-dataset-source-distribution}
\begin{tabular}{lrrrr}
\toprule
\textbf{Sorgente} & \textbf{\texttt{weapon}} & \textbf{\texttt{non\_weapon}}
                  & \textbf{\texttt{ood}} & \textbf{Totale} \\
\midrule
\texttt{01\_kaggle\_weapon}    & 100 &   0 & 100 &  200 \\
\texttt{02\_deepfirearm}       & 100 &   0 &   0 &  100 \\
\texttt{03\_google\_scraped}   & 100 &  30 &  71 &  201 \\
\texttt{04\_telegram\_youtube} & 100 &  57 & 225 &  382 \\
\texttt{05\_deepweb}           & 100 & 413 & 104 &  617 \\
\midrule
\textbf{Totale} & \textbf{500} & \textbf{500} & \textbf{500} & \textbf{1500} \\
\bottomrule
\end{tabular}
\end{table}

Il protocollo di selezione è stato condotto da un singolo revisore, secondo
criteri espliciti e documentati. Questa scelta rappresenta un limite rispetto a
protocolli multi-revisore con misura formale di accordo inter-annotatore, ma
consente di mantenere piena tracciabilità delle decisioni, coerenza applicativa
dei criteri e riproducibilità documentale del processo di freezing del dataset.

A partire dal dataset frozen viene inoltre generato il subset binario ufficiale
per gli esperimenti adversariali e anti-forensi, memorizzato in
\texttt{datasets/final/manifests/manual\_selection\_adversarial\_subset.csv}.
Questo subset contiene esclusivamente le classi \texttt{weapon} e
\texttt{non\_weapon}, per un totale di 1000 immagini bilanciate. I campioni
\gls{ood} non vengono utilizzati come target primari per la generazione delle
perturbazioni, poiché rispondono a una domanda sperimentale distinta: valutare
se modelli e strumenti automatici tendano ad assegnare impropriamente una classe
binaria nota a contenuti borderline, anomali o semanticamente esterni al dominio
operativo del task.

\begin{table}[H]
\centering
\caption{Distribuzione del subset binario per esperimenti adversariali
         e anti-forensi.}
\label{tab:methodology-adversarial-subset}
\begin{tabularx}{0.75\textwidth}{@{}l>{\raggedleft\arraybackslash}X@{}}
\toprule
\textbf{Classe} & \textbf{Numero di immagini} \\
\midrule
\texttt{weapon}      & 500 \\
\texttt{non\_weapon} & 500 \\
\midrule
\textbf{Totale}      & \textbf{1000} \\
\bottomrule
\end{tabularx}
\end{table}

La presenza di un dataset congelato rappresenta un requisito metodologico
centrale. Una volta prodotti i manifest finali, le fasi successive della
pipeline devono riferirsi a questi file come fonte ufficiale, evitando
modifiche implicite al corpus sperimentale. In questo modo, eventuali variazioni
nei risultati potranno essere attribuite ai modelli, alle perturbazioni, agli
strumenti o alle metriche, e non a cambiamenti non documentati nella
composizione del dataset.


% ─────────────────────────────────────────────────────────────────────────────
\section{Split generation and OOD evaluation set}
\label{sec:methodology-split-generation}
% ─────────────────────────────────────────────────────────────────────────────

La fase successiva della pipeline dataset consiste nella generazione degli split
puliti e del set \gls{ood}. Questa operazione è eseguita mediante lo script
\texttt{datasets/scripts/splits/11\_generate\_clean\_and\_ood\_splits.py},
che utilizza come input i due manifest ufficiali prodotti nella fase di
freezing: \texttt{manual\_selection\_final\_1500.csv} e
\texttt{manual\_selection\_adversarial\_subset.csv}. Lo script non modifica
i manifest finali, ma produce copie organizzate delle immagini e manifest
specifici per le successive fasi sperimentali.

Il subset binario composto dalle classi \texttt{weapon} e \texttt{non\_weapon}
viene suddiviso in cinque fold deterministici, denominati \texttt{fold\_1},
\texttt{fold\_2}, \texttt{fold\_3}, \texttt{fold\_4} e \texttt{fold\_5}.
Ciascun fold contiene 200 immagini, equamente distribuite tra le due classi
operative del task binario. La procedura di split utilizza un seed pseudo-casuale
fissato a 42, in modo che la partizione possa essere rigenerata in modo
riproducibile a partire dai manifest ufficiali e dallo script indicato.
In questa fase, il termine \texttt{fold} indica unità sperimentali bilanciate e
tracciabili, non partizioni di addestramento nel senso della cross-validation
tradizionale. L'impiego dei fold nelle fasi di addestramento e valutazione dei
modelli viene descritto nella Sezione~\ref{sec:methodology-models-classification}.

\begin{table}[H]
\centering
\caption{Distribuzione degli split puliti generati dal subset binario.}
\label{tab:methodology-clean-folds}
\begin{tabularx}{\textwidth}{@{}l>{\raggedleft\arraybackslash}X
                               >{\raggedleft\arraybackslash}X
                               >{\raggedleft\arraybackslash}X@{}}
\toprule
\textbf{Fold} & \textbf{\texttt{weapon}} & \textbf{\texttt{non\_weapon}}
              & \textbf{Totale} \\
\midrule
\texttt{fold\_1} & 100 & 100 & 200 \\
\texttt{fold\_2} & 100 & 100 & 200 \\
\texttt{fold\_3} & 100 & 100 & 200 \\
\texttt{fold\_4} & 100 & 100 & 200 \\
\texttt{fold\_5} & 100 & 100 & 200 \\
\midrule
\textbf{Totale} & \textbf{500} & \textbf{500} & \textbf{1000} \\
\bottomrule
\end{tabularx}
\end{table}

La classe \texttt{ood} viene invece mantenuta separata dal subset binario e
organizzata in un unico set di valutazione, collocato in
\texttt{datasets/splits/ood/ood\_eval\_set/ood/}. Questa scelta riflette la
diversa funzione metodologica dei campioni fuori distribuzione. Essi non sono
utilizzati per generare perturbazioni adversariali o anti-forensi su classi
note, ma per valutare il comportamento dei sistemi quando ricevono contenuti
che non dovrebbero essere trattati come esempi ordinari di \texttt{weapon} o
\texttt{non\_weapon}. La separazione del set \gls{ood} consente quindi di
distinguere la robustezza sul task binario dalla gestione di input borderline,
anomali o semanticamente esterni al dominio operativo.

Gli output principali della fase di split generation sono tre artefatti:
\texttt{clean\_folds\_manifest.csv}, \texttt{ood\_eval\_manifest.csv} e
\texttt{split\_generation\_summary.json}. I manifest non sono semplici elenchi
di file, ma artefatti di controllo che registrano, per ciascun campione,
l'identificativo dell'immagine, la classe finale, il fold o il set di
appartenenza, la sorgente, il percorso dell'immagine preparata, il percorso
nello split, l'hash SHA256, l'MD5 e le informazioni relative al tipo di
campione. Il file di summary documenta i controlli di consistenza eseguiti
sulla generazione degli split, inclusa l'unicità degli \texttt{image\_id},
l'unicità degli hash e la corrispondenza tra i conteggi attesi e quelli
effettivamente prodotti.

Questa organizzazione consente di disporre di unità sperimentali stabili per
le successive valutazioni baseline, per la generazione degli input adversariali
e delle trasformazioni anti-forensi e per il confronto con gli output prodotti
dagli strumenti forensi. La separazione tra clean folds e set \gls{ood} evita
inoltre di aggregare in una sola metrica problemi metodologicamente distinti:
da un lato, la stabilità della classificazione su classi note; dall'altro, la
capacità del sistema di gestire input che non appartengono alla distribuzione
operativa del task binario.





% ─────────────────────────────────────────────────────────────────────────────
\section{Models and classification paradigm}
\label{sec:methodology-models-classification}
% ─────────────────────────────────────────────────────────────────────────────

La pipeline accademica di \gls{fairlab} utilizza un insieme di proxy model
trasparenti e riproducibili per rappresentare il comportamento di classificatori
\gls{ai}-based sottoposti a dati puliti, input adversariali e trasformazioni
anti-forensi. Tali modelli non devono essere interpretati come strumenti forensi
commerciali, né come sostituti dei software analizzati nella pipeline
operativo-forense. Essi costituiscono invece componenti sperimentali controllate,
necessarie per definire un riferimento algoritmico ispezionabile, generare
perturbazioni model-dependent e studiare la possibile trasferibilità degli
effetti osservati verso altri sistemi di classificazione.

Il paradigma di classificazione adottato dai proxy model è binario. Le etichette
operative sono \texttt{non\_weapon} e \texttt{weapon}, rispettivamente codificate
come classe 0 e classe 1 nel registro dei modelli. La classe \texttt{ood}, pur
essendo parte del dataset finale congelato, non viene utilizzata per addestrare
i proxy model e non costituisce un target primario per la generazione delle
perturbazioni adversariali o anti-forensi. Essa rimane separata per la valutazione
del comportamento fuori distribuzione, in coerenza con la distinzione metodologica
introdotta nelle Sezioni~\ref{sec:methodology-research-design} e
\ref{sec:methodology-split-generation}. Questa scelta evita di trasformare il
problema \gls{ood} in una terza classe supervisionata ordinaria e consente di
mantenere separata la valutazione della robustezza binaria dalla gestione di
input borderline, anomali o semanticamente esterni al dominio operativo.

I proxy model ufficiali considerati in questa fase sono ResNet18,
EfficientNet-B0 e \gls{clip}. ResNet18 rappresenta una baseline convoluzionale
consolidata e relativamente compatta, utile per disporre di un riferimento
\gls{cnn} standard nel task binario \cite{he2016deep}. EfficientNet-B0 viene
utilizzato come modello convoluzionale più efficiente e costituisce il primary
proxy target per la generazione degli attacchi adversariali model-dependent
nella pipeline corrente \cite{tan2019efficientnet}. \gls{clip}, invece, viene
impiegato nella forma di visual encoder congelato con una testa binaria addestrata
sul task \texttt{non\_weapon}/\texttt{weapon}, così da includere nel protocollo
un modello vision-language pre-addestrato e successivamente adattato alla
classificazione binaria \cite{radford2021learning}. La Tabella~\ref{tab:methodology-proxy-models}
riassume il ruolo metodologico dei tre modelli.

\begin{table}[H]
\centering
\caption{Proxy model utilizzati nella pipeline accademica.}
\label{tab:methodology-proxy-models}
\begin{tabularx}{\textwidth}{@{}l l X@{}}
\toprule
\textbf{Modello} & \textbf{Paradigma} & \textbf{Ruolo metodologico} \\
\midrule
\texttt{resnet18}
&
\gls{cnn} supervisionata
&
Baseline convoluzionale compatta per la classificazione binaria
\texttt{non\_weapon}/\texttt{weapon}. \\
\texttt{efficientnet\_b0}
&
\gls{cnn} supervisionata
&
Proxy target principale per la generazione degli attacchi adversariali
model-dependent. \\
\texttt{clip}
&
Visual encoder congelato con testa binaria
&
Modello vision-language pre-addestrato adattato al task binario mediante
classificatore supervisionato. \\
\bottomrule
\end{tabularx}
\end{table}

L'addestramento dei proxy model segue un protocollo fold-aware. Per ciascun
fold target, il modello viene addestrato sui quattro fold rimanenti e salvato
come checkpoint specifico del fold escluso. Ad esempio, il checkpoint associato
a \texttt{fold\_1} viene addestrato utilizzando \texttt{fold\_2},
\texttt{fold\_3}, \texttt{fold\_4} e \texttt{fold\_5}. Lo stesso schema viene
ripetuto per tutti i cinque fold. Questo protocollo riduce il rischio di leakage
tra immagini usate per l'addestramento del proxy model e immagini successivamente
attaccate o valutate nello stesso fold, preservando una separazione metodologica
tra training e holdout sperimentale.

\begin{table}[H]
\centering
\caption{Protocollo fold-aware per l'addestramento dei proxy model.}
\label{tab:methodology-fold-aware-training}
\begin{tabularx}{\textwidth}{@{}lX@{}}
\toprule
\textbf{Checkpoint target} & \textbf{Fold utilizzati per l'addestramento} \\
\midrule
\texttt{fold\_1.pt} & \texttt{fold\_2}, \texttt{fold\_3}, \texttt{fold\_4}, \texttt{fold\_5} \\
\texttt{fold\_2.pt} & \texttt{fold\_1}, \texttt{fold\_3}, \texttt{fold\_4}, \texttt{fold\_5} \\
\texttt{fold\_3.pt} & \texttt{fold\_1}, \texttt{fold\_2}, \texttt{fold\_4}, \texttt{fold\_5} \\
\texttt{fold\_4.pt} & \texttt{fold\_1}, \texttt{fold\_2}, \texttt{fold\_3}, \texttt{fold\_5} \\
\texttt{fold\_5.pt} & \texttt{fold\_1}, \texttt{fold\_2}, \texttt{fold\_3}, \texttt{fold\_4} \\
\bottomrule
\end{tabularx}
\end{table}

Lo script ufficiale per l'addestramento dei proxy model è
\texttt{models/scripts/12\_train\_proxy\_models.py}. Esso utilizza come input
il manifest \texttt{datasets/splits/manifests/clean\_folds\_manifest.csv},
applica il mapping binario \texttt{non\_weapon}~$\rightarrow 0$ e
\texttt{weapon}~$\rightarrow 1$, genera i checkpoint in
\texttt{models/checkpoints/<model\_name>/<fold>.pt} e aggiorna il registro
\texttt{models/model\_registry.json}. Il registro dei modelli documenta, per
ciascuna architettura e per ciascun fold, il pattern del checkpoint, la
dimensione di input, la normalizzazione utilizzata e l'hash SHA256 del file
prodotto. La registrazione dell'hash del checkpoint è metodologicamente rilevante
perché consente di verificare che le fasi successive della pipeline, incluse la
generazione degli attacchi model-dependent e le valutazioni di robustezza, siano
riconducibili a una versione specifica e immutata del modello.

Nel protocollo sperimentale, EfficientNet-B0 assume il ruolo di primary proxy
target per la generazione degli attacchi adversariali model-dependent. Questa
scelta non implica che EfficientNet-B0 rappresenti il modello più robusto o più
accurato in senso assoluto, né che esaurisca l'analisi comparativa della tesi.
Indica invece il modello controllato rispetto al quale vengono costruite, in
questa fase della pipeline, le perturbazioni che richiedono accesso al gradiente,
al classificatore o al comportamento decisionale del target. La valutazione
successiva potrà quindi distinguere tra prestazioni del proxy target, effetti
delle perturbazioni generate rispetto a tale modello e trasferibilità verso altri
modelli o strumenti.

È infine necessario distinguere i proxy model dagli strumenti forensi
\gls{cots} analizzati nella pipeline operativo-forense. I proxy model sono
componenti accademiche aperte, configurabili e documentate, per le quali sono
noti architettura, checkpoint, mapping delle classi, protocollo di addestramento
e hash dei file. Gli strumenti forensi commerciali, al contrario, vengono trattati
come sistemi black-box o quasi black-box: l'analisi non presuppone accesso alla
loro architettura interna, ma si concentra sugli output osservabili, sulla loro
coerenza rispetto ai manifest sperimentali e sul loro comportamento in presenza
di immagini pulite, manipolate o fuori distribuzione. Questa distinzione consente
di mantenere separati il controllo sperimentale della pipeline accademica e la
valutazione operativa dei tool forensi.

% ─────────────────────────────────────────────────────────────────────────────
\section{Adversarial attack selection and generation}
\label{sec:methodology-adversarial-attacks}
% ─────────────────────────────────────────────────────────────────────────────

La pipeline adversariale di \gls{fairlab} è progettata per generare versioni
manipolate delle immagini appartenenti al subset binario
\texttt{weapon}/\texttt{non\_weapon}, preservando il collegamento tra ciascun
campione perturbato e la corrispondente immagine pulita. La generazione degli
attacchi utilizza come input il manifest
\texttt{datasets/splits/manifests/clean\_folds\_manifest.csv}, composto da
1000 immagini distribuite in cinque fold bilanciati. I campioni \texttt{ood}
non vengono utilizzati in questa fase, poiché la loro funzione metodologica è
distinta: essi servono a valutare il comportamento fuori distribuzione e non a
produrre perturbazioni adversariali rispetto a classi binarie note.

Lo script ufficiale per questa fase è
\texttt{datasets/scripts/attacks/14\_generate\_adversarial\_attacks.py}. Esso
costituisce l'entry point fold-aware per la generazione degli attacchi
adversariali e supporta sia una modalità interattiva, utile per l'esecuzione
guidata, sia una modalità a riga di comando, più adatta alla riproducibilità
sperimentale. In coerenza con il paradigma definito nella
Sezione~\ref{sec:methodology-models-classification}, lo script distingue tra
attacchi model-dependent, costruiti rispetto a un proxy model specifico, e
trasformazioni model-agnostic, applicate senza accesso al gradiente o al
comportamento decisionale del classificatore.

Gli attacchi model-dependent utilizzati nella pipeline corrente sono
\gls{fgsm}, One Pixel Attack, Sigma Zero Attack e SuperDeepFool. Essi vengono
generati rispetto a EfficientNet-B0, assunto come primary proxy target. Per ogni
immagine appartenente a un determinato fold, lo script carica il checkpoint
corrispondente da \texttt{models/checkpoints/<target\_model>/<fold>.pt}.
Questa scelta mantiene la coerenza con il protocollo fold-aware: l'immagine di
un fold viene perturbata usando il modello addestrato secondo la separazione
definita nella fase di training dei proxy model. La procedura evita quindi di
trattare il dataset come un insieme indistinto e preserva il legame tra split,
checkpoint e artefatto generato.

La trasformazione \texttt{color\_shift} viene invece trattata come perturbazione
model-agnostic. Essa modifica globalmente le componenti cromatiche
dell'immagine, senza interrogare il modello e senza utilizzare informazioni di
gradiente. Questa distinzione è metodologicamente rilevante perché consente di
separare perturbazioni costruite rispetto a un classificatore specifico da
alterazioni visive più generali, che possono simulare variazioni cromatiche o
post-processing non necessariamente ottimizzati contro un modello.

\begin{table}[H]
\centering
\caption{Attacchi adversariali generati sul subset binario.}
\label{tab:methodology-adversarial-attacks}
\begin{tabularx}{\textwidth}{@{}l l l l r@{}}
\toprule
\textbf{Attacco} & \textbf{Dipendenza} & \textbf{Target} & \textbf{Formato} & \textbf{Immagini} \\
\midrule
\texttt{fgsm}
&
model-dependent
&
\texttt{efficientnet\_b0}
&
PNG
&
1000 \\
\texttt{one\_pixel}
&
model-dependent
&
\texttt{efficientnet\_b0}
&
PNG
&
1000 \\
\texttt{sigma\_zero}
&
model-dependent
&
\texttt{efficientnet\_b0}
&
PNG
&
1000 \\
\texttt{superdeepfool}
&
model-dependent
&
\texttt{efficientnet\_b0}
&
PNG
&
1000 \\
\texttt{color\_shift}
&
model-agnostic
&
\texttt{model\_agnostic}
&
JPEG
&
1000 \\
\bottomrule
\end{tabularx}
\end{table}

La selezione degli attacchi copre famiglie diverse di manipolazione. \gls{fgsm}
rappresenta un attacco gradient-based a singolo passo, utile come baseline
adversariale efficiente e riproducibile \cite{goodfellow2015explaining}. One
Pixel Attack introduce una perturbazione sparsa, limitata a un numero minimo di
pixel modificati, e permette di valutare la sensibilità del classificatore a
variazioni localizzate \cite{su2019onepixel}. Sigma Zero Attack è incluso come
attacco orientato alla sparsità, con parametri dedicati alla regolazione della
soglia e dell'approssimazione differenziabile della norma \(L_0\)
\cite{zhao2024sigmazero}. SuperDeepFool è adottato come variante iterativa
basata su una logica di attraversamento del confine decisionale, richiamando la
famiglia di metodi ispirati a DeepFool \cite{moosavi2016deepfool}. Color Shift,
infine, non è trattato come attacco ottimizzato rispetto al modello, ma come
trasformazione adversarial-style model-agnostic, utile per introdurre una
manipolazione cromatica globale e controllata.

Per ciascuna generazione, la pipeline produce un manifest CSV e un summary JSON
dedicati. I manifest sono collocati in \texttt{attacks/manifests/} e seguono una
nomenclatura che combina famiglia adversariale, nome dell'attacco e, quando
presente, modello target. Ogni riga mantiene il collegamento con il campione
pulito originario attraverso campi quali \texttt{generated\_image\_id},
\texttt{original\_image\_id}, \texttt{fold}, \texttt{final\_label},
\texttt{clean\_relative\_path}, \texttt{perturbed\_relative\_path},
\texttt{sha256\_original}, \texttt{sha256\_perturbed}, \texttt{md5\_original}
e \texttt{md5\_perturbed}. Per gli attacchi model-dependent, il manifest
registra inoltre il modello target, il checkpoint utilizzato e l'hash SHA256 del
checkpoint, così da rendere verificabile la relazione tra immagine perturbata e
versione specifica del proxy model.

I summary JSON documentano i parametri di generazione, il manifest di input,
l'output root, il manifest prodotto, il numero di immagini attese e generate, la
distribuzione per fold, la distribuzione per classe e i controlli di integrità.
Nel caso di \gls{fgsm}, il parametro \(\epsilon\) è pari a \(8/255\), con input
size pari a 224 e output PNG. Per One Pixel Attack, il summary registra il numero
massimo di iterazioni, la dimensione della popolazione e il seed utilizzato. Per
Sigma Zero Attack sono documentati il numero di step, i parametri \(\eta\),
\(\sigma\), \(\tau\), il fattore di adattamento della soglia e la norma di
normalizzazione del gradiente. Per SuperDeepFool sono registrati il numero
massimo di iterazioni esterne, le iterazioni interne DeepFool, il numero di
projection step e la variante utilizzata. Per \texttt{color\_shift}, invece,
sono registrati formato \gls{jpeg}, qualità di compressione e parametri di
alterazione cromatica.

È importante precisare che la presente sezione descrive la selezione e la
generazione degli artefatti adversariali, non i risultati di robustezza. I campi
relativi al successo dell'attacco, quando presenti nei summary, sono trattati in
questa fase come metadati tecnici della generazione e non come evidenza
conclusiva sul comportamento dei modelli. L'analisi quantitativa degli effetti
degli attacchi, inclusa la variazione delle metriche rispetto agli input puliti,
viene demandata alla sezione dedicata al protocollo di valutazione e,
successivamente, al capitolo dei risultati sperimentali.

% ─────────────────────────────────────────────────────────────────────────────
\section{Anti-forensic transformation pipeline}
\label{sec:methodology-anti-forensic-transformations}
% ─────────────────────────────────────────────────────────────────────────────

La pipeline anti-forense di \gls{fairlab} introduce un insieme di trasformazioni
realistiche dell'artefatto immagine, applicate al subset binario
\texttt{weapon}/\texttt{non\_weapon}. A differenza degli attacchi adversariali
model-dependent descritti nella Sezione~\ref{sec:methodology-adversarial-attacks},
queste trasformazioni non sono ottimizzate rispetto al gradiente, alla funzione
di perdita o al confine decisionale di uno specifico classificatore. Esse
rappresentano invece manipolazioni plausibili del file immagine, riconducibili
a operazioni di ricompressione, ridimensionamento, sfocatura, modifica
dell'istogramma e alterazione del contrasto. In questo senso, la loro funzione
metodologica è valutare la stabilità della classificazione in presenza di
modifiche realistiche che possono alterare caratteristiche visive o statistiche
dell'immagine senza necessariamente renderla semanticamente irriconoscibile.

Lo script ufficiale per questa fase è
\texttt{datasets/scripts/attacks/13\_generate\_anti\_forensic\_attacks.py}.
Esso utilizza come input il manifest
\texttt{datasets/splits/manifests/clean\_folds\_manifest.csv} e genera gli
artefatti trasformati nella directory \texttt{attacks/anti\_forensic/}. La
struttura di output mantiene la separazione per trasformazione, fold e classe,
secondo il pattern
\texttt{attacks/anti\_forensic/<attack>/<fold>/<label>/<image\_id>\_\_<attack>.jpg}.
Tale organizzazione consente di preservare il collegamento tra ciascuna immagine
trasformata, il fold di appartenenza, la classe originaria e il corrispondente
campione pulito.

Le trasformazioni incluse nella pipeline sono \gls{jpeg} recompression,
resample~+~resize, Gaussian blur, histogram modification e contrast stretching.
La ricompressione \gls{jpeg} applica una nuova codifica lossy dell'immagine,
introducendo artefatti compatibili con scenari ordinari di salvataggio,
inoltro o conversione del file. La trasformazione resample~+~resize riduce
temporaneamente la risoluzione dell'immagine e la riporta alla dimensione
originaria mediante interpolazione, simulando degradazioni dovute a ridimensionamenti
ripetuti o a passaggi tra piattaforme. Il Gaussian blur introduce una sfocatura
controllata, utile a valutare la dipendenza del classificatore da dettagli
locali e bordi ad alta frequenza. La histogram modification modifica la
distribuzione tonale mediante equalizzazione dell'istogramma, mentre il contrast
stretching applica una normalizzazione del contrasto attraverso autocontrast
con cutoff controllato. Tali trasformazioni sono coerenti con una prospettiva
anti-forense perché intervengono sull'artefatto immagine senza introdurre,
almeno intenzionalmente, una perturbazione costruita sul modello.

\begin{table}[H]
\centering
\caption{Trasformazioni anti-forensi generate sul subset binario.}
\label{tab:methodology-anti-forensic-transformations}
\begin{tabularx}{\textwidth}{@{}lXr@{}}
\toprule
\textbf{Trasformazione} & \textbf{Operazione metodologica} & \textbf{Immagini} \\
\midrule
\texttt{jpeg\_recompression}
&
Ricompressione \gls{jpeg} con qualità controllata.
&
1000 \\
\texttt{resample\_resize}
&
Downsampling seguito da ripristino della dimensione originaria.
&
1000 \\
\texttt{gaussian\_blur}
&
Applicazione di sfocatura gaussiana con raggio controllato.
&
1000 \\
\texttt{histogram\_modification}
&
Equalizzazione dell'istogramma dell'immagine.
&
1000 \\
\texttt{contrast\_stretching}
&
Autocontrast con cutoff percentuale controllato.
&
1000 \\
\midrule
\textbf{Totale}
&
Cinque trasformazioni applicate alle 1000 immagini del subset binario.
&
\textbf{5000} \\
\bottomrule
\end{tabularx}
\end{table}

Nel run consolidato della pipeline, ciascuna trasformazione è stata applicata
alle 1000 immagini del subset binario, mantenendo la distribuzione dei cinque
fold e il bilanciamento tra le classi \texttt{weapon} e \texttt{non\_weapon}.
Per ogni trasformazione sono quindi state prodotte 200 immagini per fold, pari
a 500 immagini \texttt{weapon} e 500 immagini \texttt{non\_weapon}. Il totale
degli artefatti anti-forensi generati è pari a 5000 immagini. Questa scelta
consente di confrontare in modo simmetrico le risposte dei modelli e, nelle fasi
successive, degli strumenti forensi, rispetto a un insieme bilanciato di
manipolazioni non ottimizzate sul classificatore.

I parametri principali della generazione sono registrati nel summary
\texttt{attacks/manifests/anti\_forensic\_generation\_summary.json}. Nel run
corrente, la ricompressione \gls{jpeg} utilizza qualità 70, la trasformazione
resample~+~resize applica un fattore di scala pari a 0.5, il Gaussian blur usa
raggio 1.5 e il contrast stretching usa cutoff pari a 1.0. Tutte le immagini
generate sono salvate in formato \texttt{.jpg}; per le trasformazioni diverse
dalla ricompressione \gls{jpeg} primaria, il salvataggio avviene con qualità
\gls{jpeg} pari a 95. Questi parametri non sono presentati come ottimali, ma
come configurazione controllata e riproducibile per introdurre degradazioni
moderate e comparabili sull'intero subset sperimentale.

La pipeline produce il manifest
\texttt{attacks/manifests/anti\_forensic\_attacks\_manifest.csv}, nel quale ogni
riga descrive un'immagine trasformata e mantiene il collegamento con il campione
pulito originario. Il manifest registra campi quali
\texttt{generated\_image\_id}, \texttt{original\_image\_id}, \texttt{fold},
\texttt{final\_label}, \texttt{source\_dataset}, \texttt{clean\_relative\_path},
\texttt{perturbed\_relative\_path}, \texttt{attack\_family},
\texttt{attack\_name}, \texttt{attack\_parameters}, \texttt{sha256\_original},
\texttt{md5\_original}, \texttt{sha256\_perturbed} e
\texttt{md5\_perturbed}. La registrazione degli hash consente di verificare sia
l'identità del file pulito sia quella dell'artefatto trasformato, preservando la
tracciabilità tra input, trasformazione applicata e output generato.

Il file di sintesi documenta inoltre i controlli di consistenza della fase di
generazione: corrispondenza tra immagini attese e immagini effettivamente
prodotte, unicità degli identificativi generati, unicità degli hash SHA256 delle
immagini trasformate e corretta scrittura del manifest. Tali controlli non
costituiscono una valutazione dei risultati di robustezza, ma attestano la
coerenza tecnica della generazione degli artefatti anti-forensi. L'analisi
dell'impatto di tali trasformazioni sulle prestazioni dei modelli e, in seguito,
sugli strumenti forensi commerciali viene demandata alle sezioni dedicate al
protocollo metrico e al capitolo dei risultati.



% ─────────────────────────────────────────────────────────────────────────────
\section{Forensic evaluation bundle}
\label{sec:methodology-forensic-evaluation-bundle}
% ─────────────────────────────────────────────────────────────────────────────


% ─────────────────────────────────────────────────────────────────────────────
\section{Evaluation metrics}
\label{sec:methodology-evaluation-metrics}
% ─────────────────────────────────────────────────────────────────────────────

% ─────────────────────────────────────────────────────────────────────────────
\section{Explainability approach}
\label{sec:methodology-explainability}
% ─────────────────────────────────────────────────────────────────────────────

% ─────────────────────────────────────────────────────────────────────────────
\section{Traceability and reproducibility controls}
\label{sec:methodology-traceability-controls}
% ─────────────────────────────────────────────────────────────────────────────




\begin{comment}
% ─────────────────────────────────────────────────────────────────────────────
\section{Traceability and reproducibility controls}
\label{sec:methodology-traceability-controls1}
% ─────────────────────────────────────────────────────────────────────────────

La pipeline dataset è costruita attorno a un principio di tracciabilità progressiva. Ogni fase produce artefatti documentali che consentono di ricostruire il passaggio dalla sorgente grezza al campione incluso negli split finali. Tale impostazione è coerente con i requisiti generali della \gls{df}, nella quale l'affidabilità dell'analisi dipende anche dalla possibilità di documentare le trasformazioni subite dagli artefatti digitali e di verificarne l'identità attraverso hash e manifest \cite{casey2011digital,nist2006guide,iso27037}.

Il valore SHA256 rappresenta l'identificatore crittografico principale della pipeline. Esso viene calcolato durante la preparazione tecnica del dataset, propagato nei manifest successivi e verificato durante la copia degli artefatti negli split. L'MD5 viene inoltre registrato negli split perché alcuni strumenti forensi commerciali possono esporre o utilizzare tale valore nei report e negli hash set. Questa doppia registrazione facilita il raccordo tra gli artefatti prodotti dalla pipeline accademica e gli output esportati dagli strumenti operativi.

La Tabella~\ref{tab:methodology-dataset-artifacts} sintetizza gli artefatti principali prodotti fino allo stato attuale della pipeline dataset.

\begin{table}[H]
\centering
\caption{Principali artefatti dataset prodotti dalla pipeline.}
\label{tab:methodology-dataset-artifacts}
\begin{tabularx}{\textwidth}{@{}lX@{}}
\toprule
\textbf{Artefatto} & \textbf{Funzione metodologica} \\
\midrule
\texttt{metadata.csv} & Metadata tecnico del pool preparato, con provenienza, dimensioni, hash e informazioni di validità. \\
\texttt{review\_manifest\_full.csv} & Manifest che collega la preparazione tecnica alla revisione semantica \textit{human-in-the-loop}. \\
\texttt{manual\_selection\_protocol\_db.csv} & Database operativo della revisione manuale. \\
\texttt{manual\_selection\_final\_1500.csv} & Manifest ufficiale del dataset frozen da 1500 immagini. \\
\texttt{manual\_selection\_adversarial\_subset.csv} & Manifest ufficiale del subset binario da 1000 immagini. \\
\texttt{manual\_selection\_removed.csv} & Manifest dei campioni esclusi dalla selezione finale. \\
\texttt{clean\_folds\_manifest.csv} & Manifest degli split puliti per la valutazione baseline e la generazione delle perturbazioni. \\
\texttt{ood\_eval\_manifest.csv} & Manifest del set \gls{ood} separato. \\
\texttt{split\_generation\_summary.json} & Report dei controlli di consistenza e integrità della generazione degli split. \\
\bottomrule
\end{tabularx}
\end{table}

Questi artefatti permettono di distinguere tra dati grezzi, pool tecnico, selezione semantica, dataset congelato e split sperimentali. Tale distinzione è necessaria per preservare la riproducibilità del lavoro e per evitare ambiguità tra la fase di costruzione del dataset e le successive fasi di valutazione dei modelli. Inoltre, la disponibilità di manifest separati consente di effettuare controlli indipendenti su duplicati, percorsi mancanti, classi assegnate e distribuzione dei campioni.


% ─────────────────────────────────────────────────────────────────────────────
\section{Transition to subsequent experimental stages}
\label{sec:methodology-transition-next-stages}
% ─────────────────────────────────────────────────────────────────────────────

Gli artefatti descritti nelle sezioni precedenti costituiscono la base congelata della sperimentazione. A partire dai clean folds verranno eseguite le valutazioni baseline dei modelli \gls{ai}-based e verranno generate le versioni perturbate delle immagini mediante attacchi adversariali e trasformazioni anti-forensi. Il set \gls{ood}, invece, verrà mantenuto separato e utilizzato per valutare il comportamento dei sistemi rispetto a contenuti borderline o fuori distribuzione.

Le fasi successive prevedono la costruzione di un \textit{forensic evaluation bundle}, destinato a raccogliere in una struttura omogenea le immagini pulite, le immagini perturbate e il set \gls{ood}, insieme ai relativi manifest e hash. Tale bundle avrà la funzione di rendere trasferibili gli artefatti verso strumenti forensi commerciali come X-Ways Forensics, Magnet AXIOM, Cellebrite \gls{ufed} e Oxygen Forensics, preservando il collegamento tra file analizzato, campione originario, trasformazione applicata e output dello strumento.

Poiché la generazione delle perturbazioni e la costruzione del bundle operativo appartengono alle fasi sperimentali successive, il presente capitolo non anticipa risultati né descrive comportamenti osservati dei modelli o dei tool forensi. La metodologia qui definita stabilisce invece la base dataset e i controlli di tracciabilità necessari affinché tali valutazioni possano essere condotte in modo coerente, verificabile e riproducibile.






\end{comment}